
\scanpage{53}
\mspage{9}
\uline{NB} Ayant (via un groupoïde équivalent à description purement combinatoire) ``réconstruit''
\deleted{\unsure{du}} \deleted{\unsure{purement}} réalisé le groupoïde
$S\mathcal TD_{\mathrm{top}}^\pm$, et le foncteur
\[
 S\mathcal TD_{\mathrm{top}}^\pm\longrightarrow(\mathrm{Gr.\ pond.\ rep.})
\]
alors on aura par \unsure{cela} aussi \deleted{\unsure{une description}} \inserted{l'équivalent de}
\deleted{de} la donnée d'un $(X,\Sigma,R)$, muni \uline{en plus} d'une partie $R'\subset R$ --- \deleted{\unsure{d'où}} i.e. \unsure{avec} la construction d'un groupoïde \unsure{au-dessus}
\diagram{\begin{tikzcd}[ampersand replacement=\&,row sep=large]
S_*\mathcal TD_{\mathrm{top}}^\pm\
 \left(\begin{gathered}\text{à sous ens. marqué }R'\\
 \text{de sommets-}\\\text{repères}\end{gathered}\right)\arrow[d]\\
S_*\mathcal TD_{\mathrm{top}}^\pm
\end{tikzcd}}
et \deleted{\unsure{d'où}} il nous faudra donc \inserted{décrire} un foncteur ``gommage de $R'$''
\diagram{\begin{tikzcd}[ampersand replacement=\&,column sep=huge]
S_*\mathcal TD_{\mathrm{top}}^\pm
 \left(\underset{\text{marqué}}{\overset{\text{ens}}{R'}}\right)
 \arrow[r,"\text{gommage de }R'"]
 \&S_*\mathcal TD_{\mathrm{top}}^\pm
\end{tikzcd}}
\deleted{Cela s'exprime, sur le}
donnant lieu à un diagramme comm. de foncteurs
\diagram{\begin{tikzcd}[ampersand replacement=\&,row sep=large,column sep=huge]
S_*\mathcal TD_{\mathrm{top}}^\pm(\ldots R'\ldots)
 \arrow[r,"\text{gommage de }R'"]
 \arrow[d,"\text{oubli du }\deleted{\unsure{R'}}\text{ marquage }R'"']
\&S_*\mathcal TD_{\mathrm{top}}^\pm\arrow[dd]\\
S_*\mathcal TD\deleted{\unsure{\#}}_{\mathrm{top}}^\pm\arrow[dr]\\
\&(\mathrm{Gr}\deleted{\mathrm{bi}}\mathrm{pond})
 \arrow[r]\&
 \begin{gathered}\text{graphes avec}\\\text{pondération}\\
 \text{par }\gamma,\text{ mais}\\\text{non par }r,\end{gathered}
\end{tikzcd}}
alors qu'au niveau

\margin{\uline{NB} Cette opération a l'air d'un \uline{affaiblissement} de structure --- mais cet affaiblissement, \unsure{sauf si|si} pour une composante $C$ \unsure{donnée} de $\Sigma^+$, l'on \unsure{doive|doit} ``gommer'' \uline{tous} les repères de \deleted{\unsure{ce}} $C$ --- \unsure{alors} c'est au contraire d'un \unsure{renforcement} de la structure \unsure{qu'il s'agit}, \deleted{\unsure{plutôt}} \unsure{dans le} \unsure{renforcement} \unsure{des nouvelles données|des groupes dans ce cas} --- ou plutôt, des \unsure{groupoïdes} \unsure{du même type} $R'_C=R'\cap C$ \unsure{distincts}, et \unsure{telles} \unsure{catégories isotopiques} \unsure{de type} \unsure{(plus bas)} \unsure{avec d'éventuels} \deleted{\unsure{au dessus de}} \inserted{l'opération} de surmarquage et \unsure{l'autre} sont inverses \unsure{l'un de l'autre}, sont des équivalences de groupoïdes.
Il faudrait \unsure{distinguer les deux cas}, pour distinguer ces \uline{opérations de gommage} \unsure{essentielles|anodines} (qui \unsure{changent le groupe des} \unsure{autom.}) et les \unsure{autres}.}
\mspage{52}
